%
% daemon-lifecycle.tex
%
% Z specification of `vox daemon stop`, `vox daemon start` and
% `vox daemon restart` (bead vox-9354), modelling the outcome tables of
% docs/daemon-lifecycle-design.md section 4 and the checks of section 8.
%

\documentclass[a4paper,10pt,fleqn]{article}
\usepackage[margin=1in]{geometry}
\usepackage{fuzz}
\usepackage[colorlinks=true,linkcolor=blue,citecolor=blue,urlcolor=blue]{hyperref}

\begin{document}

\title{The voxd Lifecycle Verbs: A Z Specification}
\author{jms, for bead vox-9354}
\date{November 2026}
\hypersetup{
  pdftitle={The voxd Lifecycle Verbs: A Z Specification},
  pdfauthor={Punt Labs},
  pdfsubject={Z Specification},
  pdfcreator={fuzz/probcli}
}
\maketitle

\tableofcontents
\newpage

\section{Introduction}

The design in \texttt{docs/daemon-lifecycle-design.md} (hereafter \emph{the
design}) adds two verbs, \texttt{vox daemon stop} and \texttt{vox daemon
start}, and recomposes \texttt{restart} from them. Each verb observes the
service manager and port 8421, acts, verifies, and reports exactly one outcome
from a closed set (design \S4.4). The outcome tables of \S4.1 and \S4.2 are the
specification of the verbs; \S4.3 composes them.

This document states those tables as relations over the observation, models
the verbs as operations over a small world in which the service manager, the
processes, a foreign listener and a wheel upgrade all move independently of
vox, and lets \texttt{probcli} explore every interleaving on a set of three
process handles. Five properties are checked in every reachable state
(Section~\ref{sec:invariant}):

\begin{enumerate}
  \item[P1] every reachable observation yields exactly one row, at every stage
    of every verb (the tables are total and deterministic);
  \item[P2] inside \texttt{restart}, \texttt{start} never acts while a previous
    voxd is alive;
  \item[P3] \texttt{restart} never reports \texttt{restarted} for a serving
    process that is the previous process, is not the manager's main process,
    or runs a version other than the installed wheel (the T2-inside-restart
    path);
  \item[P4] no token is sent unless a visible voxd holds 8421, and then only
    to that process;
  \item[P5] no signal goes to a process other than a visible voxd observed on
    8421 at the escalation, and never to the manager's main process.
\end{enumerate}

A sixth, P6, checks that each verb reports only outcomes the \S4.4 table
assigns to it.

A literal transcription of the tables fails P1 in five places, and P2 fails
under one assumption about \texttt{launchd} that the design leaves open.
Wherever the design is silent or contradicts itself, the model below carries
a stated resolution, so that P1--P6 can be checked to completion.
Section~\ref{sec:findings} gives each defect with the state that triggers it,
as \texttt{probcli} found it on the literal variant.

\section{Basic and Free Types}

A Z boolean, named so that it cannot collide with the B type \texttt{BOOL}.

\begin{zed}
  ZBOOL ::= ztrue | zfalse
\end{zed}

\texttt{PID} is the bounded set of process identities. An element names a
process number, and the model allows a number to be reused by a new process
after the old one exits. Where the design holds a \emph{handle} to a process
(a \texttt{DaemonProcess}: a pidfd on Linux, a pid plus start time on macOS),
the model holds the number in a set that is pruned the moment that process
exits, so that a later process with the same number is never mistaken for it.
Three numbers suffice for every scenario in the design: a manager's process,
an unmanaged voxd, and one fresh process started while both are alive.

\begin{zed}
  PID ::= p1 | p2 | p3
\end{zed}

The three verbs, and $vNone$ for the moments between them.

\begin{zed}
  VERB ::= vNone | vStop | vStart | vRestart
\end{zed}

The result of asking the service manager to act: the command succeeded, was
refused for privilege (\texttt{PrivilegeDenied}), or failed otherwise
(\texttt{BackendCommandError}).

\begin{zed}
  RESP ::= rspOk | rspDenied | rspFailed
\end{zed}

Where a verb is in its run. A verb that is not running is $phIdle$;
$phGap$ is the moment inside \texttt{restart} after \texttt{stop} has
returned and before \texttt{start} begins.

\begin{zed}
  PHASE ::= phIdle | phStopAct | phStopWait | phStopKill | phStopReap \\
  \t1 | phGap | phStartAct | phStartWait
\end{zed}

The closed outcome set of design \S4.4, with $oNone$ for ``no report yet''.

\begin{zed}
  OUTCOME ::= oNone | oStopped | oAlreadyStopped | oStarted \\
  \t1 | oAlreadyRunning | oRestarted | oNotInstalled \\
  \t1 | oPermissionDenied | oBackendFailed | oTimeout \\
  \t1 | oPortConflict | oVersionMismatch
\end{zed}

The rows of the \texttt{stop} table (\S4.1), plus three that the table leaves
implicit: $sPE$ (the port probe failed, \S4), $sAct$ (none of S1, S2 applies,
so the verb acts) and $sGo$ (the stop request returned, or its error was
overtaken by the postcondition, so the verb waits).

\begin{zed}
  SROW ::= sPE | sS1 | sS2 | sAct | sS3 | sS4 | sGo \\
  \t1 | sS5 | sS6 | sS7 | sS8
\end{zed}

The rows of the \texttt{start} table (\S4.2). T2 splits three ways by what it
finds: $tRun$ (fresh), $tRunStale$ (stale, standalone) and $tStale$ (stale,
inside \texttt{restart}). $tCold$ is the cold start and $tGo$ the wait after
the start request.

\begin{zed}
  TROW ::= tPE | tT1 | tT4 | tT2 | tT3 | tCold | tT5 | tT6 | tGo \\
  \t1 | tT7 | tT8 | tT9 | tT10 | tRun | tRunStale | tStale
\end{zed}

\section{Constants}

\subsection{How the service manager stops}

$syncStop$ records whether the service manager's stop returns only after the
main process has exited. \texttt{systemctl stop} does: the unit is not
inactive until its control group is empty. Whether \texttt{launchctl
bootout} does is not established by the design. The model is checked with
$syncStop = ztrue$; Finding~F5 is the counterexample at $zfalse$.

\begin{axdef}
  syncStop : ZBOOL
\where
  syncStop = ztrue
\end{axdef}

\subsection{The \texttt{stop} table}

$stopPre$ classifies the observation taken before \texttt{stop} acts: whether
the probe failed, whether the unit is installed and active, which visible voxd
holds 8421, and whether a holder vox cannot see is listening. S2 is the only
row that makes no sudo call; every other installed state acts.

\begin{axdef}
  stopPre : ZBOOL \cross ZBOOL \cross ZBOOL \cross \power PID \cross ZBOOL
    \rel SROW
\where
  stopPre = \{ pe, ins, act : ZBOOL; port : \power PID; fgn : ZBOOL;
      r : SROW | \\
  \t1 (r = sPE \land pe = ztrue) \lor \\
  \t1 (r = sS1 \land pe = zfalse \land ins = zfalse) \lor \\
  \t1 (r = sS2 \land pe = zfalse \land ins = ztrue \land act = zfalse
      \land port = \emptyset \land fgn = zfalse) \lor \\
  \t1 (r = sAct \land pe = zfalse \land ins = ztrue \land \\
  \t2 \lnot (act = zfalse \land port = \emptyset \land fgn = zfalse)) \\
  \t1 @ (pe, ins, act, port, fgn) \mapsto r \}
\end{axdef}

$stopAct$ classifies the result of \texttt{request\_stop} against the
re-observation made at once after it (\S3.2): a refusal is reported only if
the unit is still active; otherwise the verb carries on to the wait.

\begin{axdef}
  stopAct : RESP \cross ZBOOL \rel SROW
\where
  stopAct = \{ rsp : RESP; a : ZBOOL; r : SROW | \\
  \t1 (r = sS3 \land rsp = rspDenied \land a = ztrue) \lor \\
  \t1 (r = sS4 \land rsp = rspFailed \land a = ztrue) \lor \\
  \t1 (r = sGo \land (rsp = rspOk \lor a = zfalse)) \\
  \t1 @ (rsp, a) \mapsto r \}
\end{axdef}

$stopSettle$ classifies the observation during the wait: whether the unit is
active, which captured handles are still alive, which visible voxd holds
8421, and whether an invisible holder listens. S5 is the confirmed stop of
\S5, conditions (a)--(d). S8 includes a captured process that has released
the port but not exited at the deadline.

\begin{axdef}
  stopSettle : ZBOOL \cross \power PID \cross \power PID \cross ZBOOL
    \rel SROW
\where
  stopSettle = \{ a : ZBOOL; cap, port : \power PID; fgn : ZBOOL;
      r : SROW | \\
  \t1 (r = sS5 \land a = zfalse \land cap = \emptyset \land port = \emptyset
      \land fgn = zfalse) \lor \\
  \t1 (r = sS6 \land a = zfalse \land port \neq \emptyset) \lor \\
  \t1 (r = sS7 \land a = zfalse \land port = \emptyset \land fgn = ztrue)
      \lor \\
  \t1 (r = sS8 \land (a = ztrue \lor \\
  \t2 (cap \neq \emptyset \land port = \emptyset \land fgn = zfalse))) \\
  \t1 @ (a, cap, port, fgn) \mapsto r \}
\end{axdef}

\subsection{The \texttt{start} table}

$startPre$ classifies the observation taken before \texttt{start} acts, in
the order the design fixes: the probe, T1, T4, then T2 or T3, then the cold
start. The last component is the set of processes answering a gated health
request; it can be non-empty only when a visible voxd holds 8421.

\begin{axdef}
  startPre : ZBOOL \cross ZBOOL \cross ZBOOL \cross \power PID \cross ZBOOL
    \cross \power PID \rel TROW
\where
  startPre = \{ pe, ins, act : ZBOOL; port : \power PID; fgn : ZBOOL;
      ans : \power PID; r : TROW | \\
  \t1 (r = tPE \land pe = ztrue) \lor \\
  \t1 (r = tT1 \land pe = zfalse \land ins = zfalse) \lor \\
  \t1 (r = tT4 \land pe = zfalse \land ins = ztrue \land act = zfalse \land
      (port \neq \emptyset \lor fgn = ztrue)) \lor \\
  \t1 (r = tT2 \land pe = zfalse \land ins = ztrue \land act = ztrue \land
      port \neq \emptyset \land ans \neq \emptyset) \lor \\
  \t1 (r = tT3 \land pe = zfalse \land ins = ztrue \land act = ztrue \land
      ans = \emptyset) \lor \\
  \t1 (r = tCold \land pe = zfalse \land ins = ztrue \land act = zfalse \land
      port = \emptyset \land fgn = zfalse) \\
  \t1 @ (pe, ins, act, port, fgn, ans) \mapsto r \}
\end{axdef}

$startAct$ classifies the result of \texttt{request\_start} against the
immediate re-observation, by the same \S3.2 rule as $stopAct$.

\begin{axdef}
  startAct : RESP \cross ZBOOL \rel TROW
\where
  startAct = \{ rsp : RESP; a : ZBOOL; r : TROW | \\
  \t1 (r = tT5 \land rsp = rspDenied \land a = zfalse) \lor \\
  \t1 (r = tT6 \land rsp = rspFailed \land a = zfalse) \lor \\
  \t1 (r = tGo \land (rsp = rspOk \lor a = ztrue)) \\
  \t1 @ (rsp, a) \mapsto r \}
\end{axdef}

$startDecide$ is the decision T2 and the cold path share once health answers
or the deadline passes. Its arguments are whether this is the cold path, the
verb (only \texttt{restart} sets \texttt{require\_fresh}), the answering
process, the manager's main process, and the processes running a version
other than the installed wheel. T8 comes before T9, as the design orders
them. T8 includes the case in which the manager reports no main process.

\begin{axdef}
  startDecide : ZBOOL \cross VERB \cross \power PID \cross \power PID
    \cross \power PID \rel TROW
\where
  startDecide = \{ cld : ZBOOL; vb : VERB; ans, mp, stl : \power PID;
      r : TROW | \\
  \t1 (r = tT7 \land ans = \emptyset) \lor \\
  \t1 (r = tT8 \land ans \neq \emptyset \land \lnot (ans \subseteq mp)) \lor \\
  \t1 (r = tT9 \land cld = ztrue \land ans \neq \emptyset \land
      ans \subseteq mp \land ans \subseteq stl) \lor \\
  \t1 (r = tT10 \land cld = ztrue \land ans \neq \emptyset \land
      ans \subseteq mp \land ans \cap stl = \emptyset) \lor \\
  \t1 (r = tRun \land cld = zfalse \land ans \neq \emptyset \land
      ans \subseteq mp \land ans \cap stl = \emptyset) \lor \\
  \t1 (r = tRunStale \land cld = zfalse \land vb \neq vRestart \land
      ans \neq \emptyset \land ans \subseteq mp \land ans \subseteq stl) \lor \\
  \t1 (r = tStale \land cld = zfalse \land vb = vRestart \land
      ans \neq \emptyset \land ans \subseteq mp \land ans \subseteq stl) \\
  \t1 @ (cld, vb, ans, mp, stl) \mapsto r \}
\end{axdef}

The outcome each decision row reports, before \texttt{restart} relabels it.

\begin{axdef}
  decideOutcome : TROW \pfun OUTCOME
\where
  decideOutcome = \{ tT7 \mapsto oTimeout, tT8 \mapsto oBackendFailed,
      tT9 \mapsto oVersionMismatch, \\
  \t2 tT10 \mapsto oStarted, tRun \mapsto oAlreadyRunning,
      tRunStale \mapsto oAlreadyRunning, \\
  \t2 tStale \mapsto oVersionMismatch \}
\end{axdef}

\section{State}
\label{sec:state}

The state has three parts. The \emph{world} is what vox observes and does not
own: whether the unit is installed and active, the manager's main process,
this user's live voxd processes, which of them holds 8421 (the visible holder,
\texttt{procs} in the design), whether a holder vox cannot see listens, which
holder answers a gated health request, and which processes run a version other
than the installed wheel. The \emph{control} is where the running verb is. The
\emph{record} is what the last step did: its outcome, the process it named,
where it sent a token, and what it signalled.

$captured$ holds the handles \texttt{stop} took at its observation, $oldMain$
the manager's main process at the same moment, and $doomed$ the handles S6
marked for termination. Each is pruned when its process exits.

The predicate is empty on purpose. In \texttt{probcli}'s reading of Z the
state predicate guards every operation, so a violation would disable an
operation silently rather than report itself. Every constraint, physical or
desired, is stated in the schema $Invariant$ of
Section~\ref{sec:invariant}, which \texttt{probcli} checks in every state.

\begin{schema}{Daemon}
  installed, active : ZBOOL \\
  mainpid, alive, onport : \power PID \\
  foreign : ZBOOL \\
  answering, stale : \power PID \\
  phase : PHASE \\
  verb : VERB \\
  cold : ZBOOL \\
  captured, oldMain, doomed : \power PID \\
  outcome : OUTCOME \\
  reported, tokenTo, signalled : \power PID
\end{schema}

\section{Initialisation}

Nothing runs and no verb is in progress. Whether the unit is installed is left
open, so both installed and uninstalled hosts are explored.

\begin{schema}{Init}
  Daemon
\where
  active = zfalse \\
  mainpid = \emptyset \land alive = \emptyset \land onport = \emptyset \\
  foreign = zfalse \\
  answering = \emptyset \land stale = \emptyset \\
  phase = phIdle \land verb = vNone \land cold = zfalse \\
  captured = \emptyset \land oldMain = \emptyset \land doomed = \emptyset \\
  outcome = oNone \\
  reported = \emptyset \land tokenTo = \emptyset \land signalled = \emptyset
\end{schema}

\section{The Environment}

Everything in this section happens without vox: a third party starts or
stops the unit, a process crashes and the manager respawns it, a foreground
voxd starts, a process binds 8421 or begins answering health, another user
takes the port, a new wheel is installed. These steps may fall between any two
steps of a verb, which is how the model explores races.

$EnvStep$ frames an environment step. It leaves the control alone, clears the
record, and prunes the three handle sets of any process that has exited. When
no verb is running it also forgets the last verb's handles.

\begin{schema}{EnvStep}
  \Delta Daemon
\where
  installed' = installed \\
  phase' = phase \land cold' = cold \\
  outcome' = oNone \\
  reported' = \emptyset \land tokenTo' = \emptyset \land
    signalled' = \emptyset \\
  phase = phIdle \implies (verb' = vNone \land captured' = \emptyset \land \\
  \t1 oldMain' = \emptyset \land doomed' = \emptyset) \\
  phase \neq phIdle \implies (verb' = verb \land
    captured' = captured \cap alive' \land \\
  \t1 oldMain' = oldMain \cap alive' \land doomed' = doomed \cap alive')
\end{schema}

A third party starts the unit; the manager launches a new process.

\begin{schema}{EnvUnitStart}
  EnvStep \\
  p? : PID
\where
  installed = ztrue \land active = zfalse \land p? \notin alive \\
  active' = ztrue \land mainpid' = \{ p? \} \land
    alive' = alive \cup \{ p? \} \\
  onport' = onport \land foreign' = foreign \\
  answering' = answering \land stale' = stale
\end{schema}

A third party stops the unit. Under $syncStop$ the main process has exited
when the stop returns; otherwise it may linger.

\begin{schema}{EnvUnitStop}
  EnvStep
\where
  active = ztrue \\
  active' = zfalse \land mainpid' = \emptyset \land foreign' = foreign \\
  (alive' = alive \setminus mainpid \land onport' = onport \setminus mainpid
    \land \\
  \t1 answering' = answering \setminus mainpid \land
    stale' = stale \setminus mainpid) \lor \\
  (syncStop = zfalse \land alive' = alive \land onport' = onport \land \\
  \t1 answering' = answering \land stale' = stale)
\end{schema}

A voxd starts outside the service manager, for example in the foreground.

\begin{schema}{EnvSpawnUnmanaged}
  EnvStep \\
  p? : PID
\where
  p? \notin alive \\
  alive' = alive \cup \{ p? \} \\
  active' = active \land mainpid' = mainpid \land onport' = onport \\
  foreign' = foreign \land answering' = answering \land stale' = stale
\end{schema}

A live voxd binds 8421, which must be free.

\begin{schema}{EnvBind}
  EnvStep \\
  p? : PID
\where
  p? \in alive \land onport = \emptyset \land foreign = zfalse \\
  onport' = \{ p? \} \\
  active' = active \land mainpid' = mainpid \land alive' = alive \\
  foreign' = foreign \land answering' = answering \land stale' = stale
\end{schema}

The holder of 8421 closes its listener but keeps running, as a voxd does
early in an orderly shutdown.

\begin{schema}{EnvUnbind}
  EnvStep
\where
  onport \neq \emptyset \\
  onport' = \emptyset \land answering' = \emptyset \\
  active' = active \land mainpid' = mainpid \land alive' = alive \\
  foreign' = foreign \land stale' = stale
\end{schema}

The holder of 8421 writes \texttt{serve.token} and starts answering health.

\begin{schema}{EnvServe}
  EnvStep
\where
  onport \neq \emptyset \land answering = \emptyset \\
  answering' = onport \\
  active' = active \land mainpid' = mainpid \land alive' = alive \\
  onport' = onport \land foreign' = foreign \land stale' = stale
\end{schema}

A process exits: it crashed, or a signal took effect. If it was the main
process the unit may stay active (\texttt{Restart=on-failure} is
re-activating it) or become inactive.

\begin{schema}{EnvExit}
  EnvStep \\
  p? : PID
\where
  p? \in alive \\
  alive' = alive \setminus \{ p? \} \land onport' = onport \setminus \{ p? \}
    \\
  answering' = answering \setminus \{ p? \} \land
    stale' = stale \setminus \{ p? \} \\
  mainpid' = mainpid \setminus \{ p? \} \land foreign' = foreign \\
  p? \notin mainpid \implies active' = active
\end{schema}

The manager respawns the main process of an active unit that has none.

\begin{schema}{EnvRespawn}
  EnvStep \\
  p? : PID
\where
  active = ztrue \land mainpid = \emptyset \land p? \notin alive \\
  mainpid' = \{ p? \} \land alive' = alive \cup \{ p? \} \\
  active' = active \land onport' = onport \\
  foreign' = foreign \land answering' = answering \land stale' = stale
\end{schema}

A process vox cannot see, such as another user's, takes 8421, and later
releases it.

\begin{schema}{EnvForeignBind}
  EnvStep
\where
  onport = \emptyset \land foreign = zfalse \\
  foreign' = ztrue \\
  active' = active \land mainpid' = mainpid \land alive' = alive \\
  onport' = onport \land answering' = answering \land stale' = stale
\end{schema}

\begin{schema}{EnvForeignRelease}
  EnvStep
\where
  foreign = ztrue \\
  foreign' = zfalse \\
  active' = active \land mainpid' = mainpid \land alive' = alive \\
  onport' = onport \land answering' = answering \land stale' = stale
\end{schema}

A new wheel is installed: every running process now runs an old version.

\begin{schema}{EnvUpgrade}
  EnvStep
\where
  stale \neq alive \\
  stale' = alive \\
  active' = active \land mainpid' = mainpid \land alive' = alive \\
  onport' = onport \land foreign' = foreign \land answering' = answering
\end{schema}

\section{Framing the Verbs}

A verb step that does not touch the world.

\begin{schema}{Hold}
  \Delta Daemon
\where
  installed' = installed \land active' = active \\
  mainpid' = mainpid \land alive' = alive \land onport' = onport \\
  foreign' = foreign \land answering' = answering \land stale' = stale
\end{schema}

The service manager's response to a stop request. A successful request makes
the unit inactive and, under $syncStop$, removes its main process. A failed
request leaves the world as it was, unless another actor brought the unit down
in the same instant; that race is what the \S3.2 re-observation exists for.

\begin{schema}{StopReaction}
  \Delta Daemon \\
  rsp? : RESP
\where
  installed' = installed \land foreign' = foreign \\
  (active' = zfalse \land mainpid' = \emptyset \land \\
  \t1 ((alive' = alive \setminus mainpid \land
    onport' = onport \setminus mainpid \land \\
  \t2 answering' = answering \setminus mainpid \land
    stale' = stale \setminus mainpid) \lor \\
  \t1 (syncStop = zfalse \land alive' = alive \land onport' = onport \land \\
  \t2 answering' = answering \land stale' = stale))) \lor \\
  (rsp? \neq rspOk \land active' = active \land mainpid' = mainpid \land
    alive' = alive \land \\
  \t1 onport' = onport \land answering' = answering \land stale' = stale)
\end{schema}

The service manager's response to a start request. A successful request makes
the unit active with a new main process, or active with none yet (the process
failed at once and the manager is re-activating it). A failed request leaves
the world alone unless another actor started the unit in the same instant.

\begin{schema}{StartReaction}
  \Delta Daemon \\
  rsp? : RESP; q? : PID
\where
  installed' = installed \land onport' = onport \land foreign' = foreign \\
  answering' = answering \land stale' = stale \\
  (q? \notin alive \land active' = ztrue \land mainpid' = \{ q? \} \land
    alive' = alive \cup \{ q? \}) \lor \\
  (active' = ztrue \land mainpid' = \emptyset \land alive' = alive) \lor \\
  (rsp? \neq rspOk \land active' = active \land mainpid' = mainpid \land
    alive' = alive)
\end{schema}

$StopEnd$ ends the \texttt{stop} half. Inside \texttt{restart}, a stop that
reports $oStopped$ or $oAlreadyStopped$ passes control to \texttt{start}
(design \S4.3); any other stop outcome is the \texttt{restart} outcome, and
\texttt{start} never runs.

\begin{schema}{StopEnd}
  \Delta Daemon \\
  res? : OUTCOME
\where
  (verb' = vRestart \land res? \in \{ oStopped, oAlreadyStopped \} \land
    phase' = phGap \land outcome' = oNone) \lor \\
  (\lnot (verb' = vRestart \land res? \in \{ oStopped, oAlreadyStopped \})
    \land phase' = phIdle \land outcome' = res?)
\end{schema}

$StartEnd$ ends \texttt{start}. Inside \texttt{restart} it applies the \S4.3
relabelling: $oStarted$ becomes $oRestarted$; $oAlreadyRunning$ becomes
$oRestarted$ when the serving process is not a captured previous process, and
$oBackendFailed$ (``stop did not take effect'') when it is. The comparison is
by handle, not by process number; Finding~F4 explains why.

\begin{schema}{StartEnd}
  \Delta Daemon \\
  res? : OUTCOME
\where
  phase' = phIdle \\
  (verb' = vRestart \land (res? = oStarted \lor \\
  \t1 (res? = oAlreadyRunning \land reported' \cap captured = \emptyset))
    \land outcome' = oRestarted) \lor \\
  (verb' = vRestart \land res? = oAlreadyRunning \land
    reported' \cap captured \neq \emptyset \land
    outcome' = oBackendFailed) \lor \\
  (\lnot (verb' = vRestart \land res? \in \{ oStarted, oAlreadyRunning \})
    \land outcome' = res?)
\end{schema}

\section{Stop}

Each operation below is one row group of the \S4.1 table. The verb is entered
with $vb? = vStop$ for \texttt{stop} and $vb? = vRestart$ for the first half of
\texttt{restart}; $pe? = ztrue$ means the port probe raised
\texttt{PortProbeError}.

\subsection{Before acting: probe failure, S1, S2}

\begin{schema}{StopProbeFailed}
  Hold \\
  StopEnd \\
  vb?: VERB; pe?: ZBOOL
\where
  phase = phIdle \land vb? \in \{ vStop, vRestart \} \\
  ((pe?, installed, active, onport, foreign) \mapsto sPE) \in stopPre \\
  res? = oBackendFailed \land verb' = vb? \land cold' = cold \\
  captured' = \emptyset \land oldMain' = \emptyset \land doomed' = \emptyset \\
  reported' = \emptyset \land tokenTo' = \emptyset \land
    signalled' = \emptyset
\end{schema}

\begin{schema}{StopNotInstalled}
  Hold \\
  StopEnd \\
  vb?: VERB; pe?: ZBOOL
\where
  phase = phIdle \land vb? \in \{ vStop, vRestart \} \\
  ((pe?, installed, active, onport, foreign) \mapsto sS1) \in stopPre \\
  res? = oNotInstalled \land verb' = vb? \land cold' = cold \\
  captured' = \emptyset \land oldMain' = \emptyset \land doomed' = \emptyset \\
  reported' = \emptyset \land tokenTo' = \emptyset \land
    signalled' = \emptyset
\end{schema}

\begin{schema}{StopAlreadyStopped}
  Hold \\
  StopEnd \\
  vb?: VERB; pe?: ZBOOL
\where
  phase = phIdle \land vb? \in \{ vStop, vRestart \} \\
  ((pe?, installed, active, onport, foreign) \mapsto sS2) \in stopPre \\
  res? = oAlreadyStopped \land verb' = vb? \land cold' = cold \\
  captured' = \emptyset \land oldMain' = \emptyset \land doomed' = \emptyset \\
  reported' = \emptyset \land tokenTo' = \emptyset \land
    signalled' = \emptyset
\end{schema}

Otherwise \texttt{stop} captures a handle on each visible voxd holding 8421,
notes the manager's main process, and goes on to act.

\begin{schema}{StopObserve}
  Hold \\
  vb?: VERB; pe?: ZBOOL
\where
  phase = phIdle \land vb? \in \{ vStop, vRestart \} \\
  ((pe?, installed, active, onport, foreign) \mapsto sAct) \in stopPre \\
  phase' = phStopAct \land verb' = vb? \land cold' = cold \\
  captured' = onport \land oldMain' = mainpid \land doomed' = \emptyset \\
  outcome' = oNone \\
  reported' = \emptyset \land tokenTo' = \emptyset \land
    signalled' = \emptyset
\end{schema}

\subsection{Acting: S3, S4, or wait}

\begin{schema}{StopDenied}
  StopReaction \\
  StopEnd
\where
  phase = phStopAct \\
  ((rsp?, active') \mapsto sS3) \in stopAct \\
  res? = oPermissionDenied \land verb' = verb \land cold' = cold \\
  captured' = captured \cap alive' \land oldMain' = oldMain \cap alive' \land
    doomed' = doomed \\
  reported' = \emptyset \land tokenTo' = \emptyset \land
    signalled' = \emptyset
\end{schema}

\begin{schema}{StopCommandFailed}
  StopReaction \\
  StopEnd
\where
  phase = phStopAct \\
  ((rsp?, active') \mapsto sS4) \in stopAct \\
  res? = oBackendFailed \land verb' = verb \land cold' = cold \\
  captured' = captured \cap alive' \land oldMain' = oldMain \cap alive' \land
    doomed' = doomed \\
  reported' = \emptyset \land tokenTo' = \emptyset \land
    signalled' = \emptyset
\end{schema}

\begin{schema}{StopRequested}
  StopReaction
\where
  phase = phStopAct \\
  ((rsp?, active') \mapsto sGo) \in stopAct \\
  phase' = phStopWait \land verb' = verb \land cold' = cold \\
  captured' = captured \cap alive' \land oldMain' = oldMain \cap alive' \land
    doomed' = doomed \\
  outcome' = oNone \\
  reported' = \emptyset \land tokenTo' = \emptyset \land
    signalled' = \emptyset
\end{schema}

\subsection{Waiting: S5, S6, S7, S8}

S5 fires the moment its condition holds; the others are what the poll sees at
the deadline. Because environment steps may come first, \texttt{probcli}
explores every moment at which the deadline can fall.

\begin{schema}{StopConfirmed}
  Hold \\
  StopEnd
\where
  phase = phStopWait \\
  ((active, captured, onport, foreign) \mapsto sS5) \in stopSettle \\
  res? = oStopped \land verb' = verb \land cold' = cold \\
  captured' = captured \land oldMain' = oldMain \land doomed' = doomed \\
  reported' = \emptyset \land tokenTo' = \emptyset \land
    signalled' = \emptyset
\end{schema}

S6 marks every visible voxd on 8421 for termination.

\begin{schema}{StopEscalate}
  Hold
\where
  phase = phStopWait \\
  ((active, captured, onport, foreign) \mapsto sS6) \in stopSettle \\
  phase' = phStopKill \land verb' = verb \land cold' = cold \\
  captured' = captured \land oldMain' = oldMain \land doomed' = onport \\
  outcome' = oNone \\
  reported' = \emptyset \land tokenTo' = \emptyset \land
    signalled' = \emptyset
\end{schema}

\begin{schema}{StopHolderTimeout}
  Hold \\
  StopEnd
\where
  phase = phStopWait \\
  ((active, captured, onport, foreign) \mapsto sS7) \in stopSettle \\
  res? = oTimeout \land verb' = verb \land cold' = cold \\
  captured' = captured \land oldMain' = oldMain \land doomed' = doomed \\
  reported' = \emptyset \land tokenTo' = \emptyset \land
    signalled' = \emptyset
\end{schema}

\begin{schema}{StopTimeout}
  Hold \\
  StopEnd
\where
  phase = phStopWait \\
  ((active, captured, onport, foreign) \mapsto sS8) \in stopSettle \\
  res? = oTimeout \land verb' = verb \land cold' = cold \\
  captured' = captured \land oldMain' = oldMain \land doomed' = doomed \\
  reported' = \emptyset \land tokenTo' = \emptyset \land
    signalled' = \emptyset
\end{schema}

\subsection{Escalation: terminate, then confirm S5}

$StopKill$ signals every marked handle that is still alive. The signals take
effect, or fail to, through $EnvExit$; then the verb confirms S5 or times out.

\begin{schema}{StopKill}
  Hold
\where
  phase = phStopKill \\
  phase' = phStopReap \land verb' = verb \land cold' = cold \\
  captured' = captured \land oldMain' = oldMain \land doomed' = doomed \\
  outcome' = oNone \\
  reported' = \emptyset \land tokenTo' = \emptyset \land
    signalled' = doomed
\end{schema}

\begin{schema}{StopKillConfirmed}
  Hold \\
  StopEnd
\where
  phase = phStopReap \\
  ((active, captured, onport, foreign) \mapsto sS5) \in stopSettle \\
  res? = oStopped \land verb' = verb \land cold' = cold \\
  captured' = captured \land oldMain' = oldMain \land doomed' = doomed \\
  reported' = \emptyset \land tokenTo' = \emptyset \land
    signalled' = \emptyset
\end{schema}

\begin{schema}{StopKillTimeout}
  Hold \\
  StopEnd
\where
  phase = phStopReap \\
  ((active, captured, onport, foreign) \mapsto sS5) \notin stopSettle \\
  res? = oTimeout \land verb' = verb \land cold' = cold \\
  captured' = captured \land oldMain' = oldMain \land doomed' = doomed \\
  reported' = \emptyset \land tokenTo' = \emptyset \land
    signalled' = \emptyset
\end{schema}

\section{Start}

\texttt{start} is entered from $phIdle$ as itself, or from $phGap$ as the
second half of \texttt{restart}, which keeps $verb = vRestart$ and so sets
\texttt{require\_fresh}.

\begin{schema}{StartEntry}
  \Delta Daemon
\where
  (phase = phIdle \land verb' = vStart) \lor
    (phase = phGap \land verb' = verb)
\end{schema}

\subsection{Before acting: probe failure, T1, T4, T2, T3}

\begin{schema}{StartProbeFailed}
  Hold \\
  StartEntry \\
  StartEnd \\
  pe?: ZBOOL
\where
  ((pe?, installed, active, onport, foreign, answering) \mapsto tPE)
    \in startPre \\
  res? = oBackendFailed \land cold' = cold \\
  captured' = captured \land oldMain' = oldMain \land doomed' = doomed \\
  reported' = \emptyset \land tokenTo' = \emptyset \land
    signalled' = \emptyset
\end{schema}

\begin{schema}{StartNotInstalled}
  Hold \\
  StartEntry \\
  StartEnd \\
  pe?: ZBOOL
\where
  ((pe?, installed, active, onport, foreign, answering) \mapsto tT1)
    \in startPre \\
  res? = oNotInstalled \land cold' = cold \\
  captured' = captured \land oldMain' = oldMain \land doomed' = doomed \\
  reported' = \emptyset \land tokenTo' = \emptyset \land
    signalled' = \emptyset
\end{schema}

T4 reports the holder before any token goes out.

\begin{schema}{StartPortConflict}
  Hold \\
  StartEntry \\
  StartEnd \\
  pe?: ZBOOL
\where
  ((pe?, installed, active, onport, foreign, answering) \mapsto tT4)
    \in startPre \\
  res? = oPortConflict \land cold' = cold \\
  captured' = captured \land oldMain' = oldMain \land doomed' = doomed \\
  reported' = \emptyset \land tokenTo' = \emptyset \land
    signalled' = \emptyset
\end{schema}

T2: a visible voxd holds 8421 and answers the gated health request. The token
goes to that process; the decision is $startDecide$ on the warm path.

\begin{schema}{StartRunning}
  Hold \\
  StartEntry \\
  StartEnd \\
  pe?: ZBOOL
\where
  ((pe?, installed, active, onport, foreign, answering) \mapsto tT2)
    \in startPre \\
  \exists d : TROW | ((zfalse, verb', answering, mainpid, stale) \mapsto d)
    \in startDecide @ \\
  \t1 d \in \dom decideOutcome \land res? = decideOutcome~d \\
  cold' = zfalse \\
  captured' = captured \land oldMain' = oldMain \land doomed' = doomed \\
  reported' = answering \land tokenTo' = onport \land
    signalled' = \emptyset
\end{schema}

T3: active but not yet answering. The verb waits; any token the gated probe
sends goes only to a visible holder.

\begin{schema}{StartAwait}
  Hold \\
  StartEntry \\
  pe?: ZBOOL
\where
  ((pe?, installed, active, onport, foreign, answering) \mapsto tT3)
    \in startPre \\
  phase' = phStartWait \land cold' = zfalse \\
  captured' = captured \land oldMain' = oldMain \land doomed' = doomed \\
  outcome' = oNone \\
  reported' = \emptyset \land tokenTo' = onport \land
    signalled' = \emptyset
\end{schema}

The cold start: installed, inactive, and nothing on 8421.

\begin{schema}{StartObserve}
  Hold \\
  StartEntry \\
  pe?: ZBOOL
\where
  ((pe?, installed, active, onport, foreign, answering) \mapsto tCold)
    \in startPre \\
  phase' = phStartAct \land cold' = ztrue \\
  captured' = captured \land oldMain' = oldMain \land doomed' = doomed \\
  outcome' = oNone \\
  reported' = \emptyset \land tokenTo' = \emptyset \land
    signalled' = \emptyset
\end{schema}

\subsection{Acting: T5, T6, or wait}

\begin{schema}{StartDenied}
  StartReaction \\
  StartEnd
\where
  phase = phStartAct \\
  ((rsp?, active') \mapsto tT5) \in startAct \\
  res? = oPermissionDenied \land verb' = verb \land cold' = cold \\
  captured' = captured \cap alive' \land oldMain' = oldMain \cap alive' \land
    doomed' = doomed \\
  reported' = \emptyset \land tokenTo' = \emptyset \land
    signalled' = \emptyset
\end{schema}

\begin{schema}{StartCommandFailed}
  StartReaction \\
  StartEnd
\where
  phase = phStartAct \\
  ((rsp?, active') \mapsto tT6) \in startAct \\
  res? = oBackendFailed \land verb' = verb \land cold' = cold \\
  captured' = captured \cap alive' \land oldMain' = oldMain \cap alive' \land
    doomed' = doomed \\
  reported' = \emptyset \land tokenTo' = \emptyset \land
    signalled' = \emptyset
\end{schema}

\begin{schema}{StartRequested}
  StartReaction
\where
  phase = phStartAct \\
  ((rsp?, active') \mapsto tGo) \in startAct \\
  phase' = phStartWait \land verb' = verb \land cold' = cold \\
  captured' = captured \cap alive' \land oldMain' = oldMain \cap alive' \land
    doomed' = doomed \\
  outcome' = oNone \\
  reported' = \emptyset \land tokenTo' = \emptyset \land
    signalled' = \emptyset
\end{schema}

\subsection{Waiting: T7, T8, T9, T10, or T2 after T3}

The poll returns the moment health answers, and decides; at the deadline with
no answer it reports T7.

\begin{schema}{StartAnswered}
  Hold \\
  StartEnd
\where
  phase = phStartWait \land answering \neq \emptyset \\
  \exists d : TROW | ((cold, verb, answering, mainpid, stale) \mapsto d)
    \in startDecide @ \\
  \t1 d \in \dom decideOutcome \land res? = decideOutcome~d \\
  verb' = verb \land cold' = cold \\
  captured' = captured \land oldMain' = oldMain \land doomed' = doomed \\
  reported' = answering \land tokenTo' = onport \land
    signalled' = \emptyset
\end{schema}

\begin{schema}{StartTimeout}
  Hold \\
  StartEnd
\where
  phase = phStartWait \\
  ((cold, verb, answering, mainpid, stale) \mapsto tT7) \in startDecide \\
  res? = oTimeout \land verb' = verb \land cold' = cold \\
  captured' = captured \land oldMain' = oldMain \land doomed' = doomed \\
  reported' = \emptyset \land tokenTo' = onport \land
    signalled' = \emptyset
\end{schema}

\section{Restart}

There is no \texttt{Restart} operation of its own, by design (\S4.3:
``\texttt{restart} adds no stop or start logic of its own''). It is
\[
  Restart \defs Stop[vb? = vRestart] \semi Env^{*} \semi Start
\]
in which $Stop$ is any path through the operations of the \texttt{stop}
section entered with $vb? = vRestart$, $StopEnd$ admits $phGap$ only after
$oStopped$ or $oAlreadyStopped$, $Env^{*}$ is any sequence of environment
steps (a third party may start the unit in the gap), and $Start$ is any path
through the \texttt{start} section entered from $phGap$, where
$verb = vRestart$ sets \texttt{require\_fresh} in $startDecide$ and
$StartEnd$ applies the relabelling.

\section{The Invariant}
\label{sec:invariant}

\texttt{probcli} checks $Invariant$ in every reachable state. Its first group
is the physics of the world; the remaining groups are P1--P6.

\begin{schema}{Invariant}
  Daemon
\where
  \# mainpid \leq 1 \land mainpid \subseteq alive \\
  mainpid \neq \emptyset \implies active = ztrue \\
  active = ztrue \implies installed = ztrue \\
  onport \subseteq alive \land \# onport \leq 1 \\
  foreign = ztrue \implies onport = \emptyset \\
  answering \subseteq onport \land stale \subseteq alive \\
  captured \subseteq alive \land oldMain \subseteq alive \land
    doomed \subseteq alive \\
  \forall pe : ZBOOL @ \\
  \t1 (phase = phIdle \implies \\
  \t2 \# \{ r : SROW | ((pe, installed, active, onport, foreign) \mapsto r)
    \in stopPre \} = 1) \land \\
  \t1 ((phase = phIdle \lor phase = phGap) \implies \\
  \t2 \# \{ r : TROW | ((pe, installed, active, onport, foreign, answering)
    \mapsto r) \in startPre \} = 1) \\
  (phase = phIdle \lor phase = phGap) \land active = ztrue \land
    answering \neq \emptyset \implies \\
  \t1 (\forall vb : VERB | vb \in \{ vStart, vRestart \} @ \\
  \t2 \# \{ r : TROW | ((zfalse, vb, answering, mainpid, stale) \mapsto r)
    \in startDecide \} = 1) \\
  phase = phStopAct \implies (\forall rsp : RESP; a : ZBOOL @ \\
  \t1 \# \{ r : SROW | ((rsp, a) \mapsto r) \in stopAct \} = 1) \\
  phase \in \{ phStopWait, phStopReap \} \implies \\
  \t1 \# \{ r : SROW | ((active, captured, onport, foreign) \mapsto r)
    \in stopSettle \} = 1 \\
  phase = phStartAct \implies (\forall rsp : RESP; a : ZBOOL @ \\
  \t1 \# \{ r : TROW | ((rsp, a) \mapsto r) \in startAct \} = 1) \\
  phase = phStartWait \implies \\
  \t1 \# \{ r : TROW | ((cold, verb, answering, mainpid, stale) \mapsto r)
    \in startDecide \} = 1 \\
  verb = vRestart \land phase \in \{ phGap, phStartAct, phStartWait \}
    \implies \\
  \t1 captured = \emptyset \land oldMain = \emptyset \\
  verb = vRestart \land outcome = oRestarted \implies \\
  \t1 reported \neq \emptyset \land reported \subseteq mainpid \land
    reported \cap stale = \emptyset \land \\
  \t1 reported \cap (captured \cup oldMain) = \emptyset \\
  tokenTo \subseteq onport \\
  signalled \subseteq doomed \land doomed \cap mainpid = \emptyset \\
  verb = vStop \implies outcome \in \{ oNone, oStopped, oAlreadyStopped,
    oNotInstalled, \\
  \t1 oPermissionDenied, oBackendFailed, oTimeout \} \\
  verb = vStart \implies outcome \in \{ oNone, oStarted, oAlreadyRunning,
    oNotInstalled, \\
  \t1 oPermissionDenied, oBackendFailed, oTimeout, oPortConflict,
    oVersionMismatch \} \\
  verb = vRestart \implies outcome \in \{ oNone, oRestarted, oNotInstalled,
    oPermissionDenied, \\
  \t1 oBackendFailed, oTimeout, oPortConflict, oVersionMismatch \}
\end{schema}

The clauses read as follows.

\begin{description}
  \item[Physics.] At most one main process, alive, and only while the unit is
    active; an active unit is installed; at most one visible holder of 8421,
    alive, and none while an invisible holder listens; only the holder
    answers; the handle sets name live processes only.
  \item[P1.] At each stage of each verb, the observation, together with every
    response the manager might give and every re-observation it might
    produce, falls in exactly one row. At $phIdle$ both verbs' entry tables
    are checked, and when a warm process answers, the T2 decision is checked
    for \texttt{start} and for \texttt{restart}.
  \item[P2.] Between the two halves of \texttt{restart}, and while
    \texttt{start} acts and waits inside it, no captured process is alive,
    and neither is the manager's main process as \texttt{stop} observed it.
  \item[P3.] A \texttt{restarted} report names a process that is the
    manager's main process, runs the installed wheel, and is neither a
    captured process nor the old main process.
  \item[P4.] Any token went to the visible voxd holding 8421. With no visible
    holder, no token went anywhere.
  \item[P5.] A signal goes only to a handle that S6 marked when it saw it
    holding 8421; no marked handle is the manager's main process. Because
    the handle sets are pruned on exit, a reused process number is never
    signalled.
  \item[P6.] Each verb reports only outcomes the \S4.4 table assigns to it;
    in particular \texttt{restart} never reports \texttt{stopped},
    \texttt{already\_stopped}, \texttt{started} or \texttt{already\_running}.
\end{description}

Two further facts follow without a separate clause. P2 empties $captured$
throughout \texttt{start} inside \texttt{restart}, so the second disjunct of
$StartEnd$, ``stop did not take effect'', can never fire. The \S4.3 branch
that reports it is dead code once the comparison is by handle (Finding~F4).
P5 holds by construction: the only operation that signals is $StopKill$, and
its target set is exactly $doomed$. The check confirms that no other
operation assigns $signalled$, and that the handles S6 marks are never the
manager's.

\section{Model Checking}
\label{sec:modelcheck}

\texttt{fuzz -t} accepts the specification with no errors. \texttt{probcli}
explores it exhaustively over the three process handles:

\begin{verbatim}
probcli docs/daemon-lifecycle.tex -model_check -nodead \
        -p MAX_OPERATIONS 500 -p TIME_OUT 60000

ALL OPERATIONS COVERED
Model checking finished, all open states visited
States analysed: 181105
Transitions fired: 1974748
No counter example found. ALL states visited.
\end{verbatim}

All 36 operations are covered. Deadlock checking is off (\texttt{-nodead}):
the environment is always able to move, so a deadlock report would say
nothing. A verb that sticks would show up as a P1 violation instead, because
every stage that waits requires exactly one enabled row.

The checks find violations when they should. Every mutation below was
caught. Making a stale warm daemon count as running inside \texttt{restart}
breaks P3. Reverting any one of the resolutions of F0--F3 in
Section~\ref{sec:findings} breaks P1. Setting $syncStop = zfalse$ breaks
P2. Comparing by process number, as in F4, breaks a clause that says
``stop did not take effect'' is reported only while a captured process
lives.

The model is small on purpose. It does not model time beyond ``the deadline
may fall at any moment'', the macOS start-time race (\S3.4 states its
residual window), the content of report details, the root and platform
preconditions (they raise before any observation), or which port a voxd off
8421 is listening on.

\section{Findings from Formalisation}
\label{sec:findings}

Each finding was produced by a variant of this file in which the stated
resolution was reverted to the design's literal wording, then model-checked.
The triggering state is the one \texttt{probcli} reported. None of them is
fixed by editing the design here; that is for the design's owner.

\begin{enumerate}
  \item[\textbf{F0.}] \textbf{Probe failure and not-installed overlap.}
    With S1 and T1 read literally (``not installed'') next to \S4's rule
    that a probe failure returns \texttt{backend\_failed} before acting,
    the initial state of an uninstalled host whose probe raises matches two
    rows: $stopPre$ maps $(ztrue, zfalse, zfalse, \emptyset, zfalse)$ to
    both $sPE$ and $sS1$, and $startPre$ likewise to $tPE$ and $tT1$.
    Neither row acts, so the order is not settled by ``before acting''. S1's
    detail (``names any voxd on 8421'') needs the probe that just failed.
    \emph{Model:} the probe is checked first. \emph{Recommend:} say so in
    \S4.

  \item[\textbf{F1.}] \textbf{S3 and T5 omit the \S3.2 re-observation.}
    \S3.2 says that after \emph{any} non-zero exit the lifecycle observes
    again and carries on if the postcondition holds. S4 and T6 say this; S3
    and T5 do not. Literal S3: unit active with main $p3$,
    \texttt{stop} reaches $phStopAct$; \texttt{request\_stop} raises
    \texttt{PrivilegeDenied} while another actor stops the unit, so the
    re-observation shows it inactive. $stopAct$ maps
    $(rspDenied, zfalse)$ to both S3 (\texttt{permission\_denied}) and the
    wait (which ends in S5, \texttt{stopped}). Literal T5 is the mirror
    image: a cold \texttt{start} whose request is refused while a third
    party starts the unit maps $(rspDenied, ztrue)$ to both T5 and the
    wait. The case is not exotic. When the unit is already inactive and only
    an unmanaged voxd holds 8421 (the S6 path), the literal S3 reports
    \texttt{permission\_denied} for a sudo refusal on a stop that had
    nothing to stop. \emph{Model:} S3 requires the re-observation to show
    the unit still active, and T5 requires it to show it still inactive.
    \emph{Recommend:} write the re-observation into S3 and T5, or say
    explicitly that \texttt{request\_stop} is skipped when the unit is
    already inactive.

  \item[\textbf{F2.}] \textbf{\texttt{stop} has no row for a captured
    process that has released 8421 but not exited.} An unmanaged voxd $p1$
    holds 8421 on an installed, inactive unit. \texttt{stop} captures
    $p1$; $p1$ begins an orderly shutdown and closes its listener but is
    still running; \texttt{request\_stop} succeeds. The wait now sees
    $(active, captured, onport, foreign) =
    (zfalse, \{p1\}, \emptyset, zfalse)$. S5 fails on condition (c), S6
    needs a voxd \emph{on 8421}, S7 needs a listener, and S8 needs the unit
    active or a process that survived \texttt{terminate()}. At the deadline
    no row applies. \emph{Model:} S8, \texttt{timeout} naming the process.
    \emph{Recommend:} either that, or extend S6 to terminate any captured
    handle that has not exited. It is already a pinned voxd, so the extension
    widens nothing, and the verb would then succeed instead of timing out.

  \item[\textbf{F3.}] \textbf{T8's ``when the backend reports a pid''
    leaves a hole.} The unit is active, its main process $p2$ has crashed,
    and \texttt{Restart=on-failure} is re-activating it, so \texttt{MainPID}
    is 0. An unmanaged voxd $p3$, on the installed wheel, binds 8421 and
    answers. \texttt{start} takes T3, waits, and $p3$ answers:
    $startDecide$ sees $(zfalse, vStart, \{p3\}, \emptyset, \emptyset)$.
    T8 is excluded (no main pid), T9 is excluded (the version is current),
    and T2-fresh and T10 both need $serving.pid = mainpid$. No row applies.
    On the cold path the same hole opens if the manager's process dies
    before the answer. \emph{Model:} T8 covers $mainpid = \emptyset$
    (\texttt{backend\_failed}: ``port 8421 answered by pid X; the service
    manager reports no main process''). \emph{Recommend:} the same; a
    process the manager does not own must not count as started, least of
    all inside \texttt{restart}.

  \item[\textbf{F4.}] \textbf{\S4.3 compares bare pids, against \S3.4.}
    The pseudocode tests \texttt{t.pid != s.previous\_pid}. The unit runs
    $p1$; \texttt{restart} stops it and S5 confirms that $p1$'s handle has
    exited. In the gap a third party starts the unit, and the kernel reuses
    the number $p1$ for the new main process, which binds, answers, and runs
    the installed wheel. \texttt{start} returns \texttt{already\_running}
    with $pid = previous\_pid$, and \texttt{restart} reports
    \texttt{backend\_failed}, ``stop did not take effect'', about a stop
    that did take effect. \emph{Model:} the comparison is by handle. P2 then
    shows the branch unreachable (Section~\ref{sec:invariant}).
    \emph{Recommend:} delete the branch. After \texttt{stopped} or
    \texttt{already\_stopped} every captured handle has exited, so a fresh
    \texttt{already\_running} inside \texttt{restart} is always a third
    party's start and always \texttt{restarted}.

  \item[\textbf{F5.}] \textbf{\S5's ``no new process is launched while the
    old one is alive'' rests on an unstated property of the backend.}
    Condition (c) covers only the handles captured from 8421. The manager's
    main process is not captured if it has not yet bound the port (the unit
    is activating). If \texttt{is\_active()} can turn false before that
    process exits ($syncStop = zfalse$), the following is reachable. The unit
    is active with main $p3$ not on 8421. \texttt{restart} observes
    ($captured = \emptyset$, $oldMain = \{p3\}$). The unit goes inactive while
    $p3$ lingers. S5 holds, since nothing is captured and 8421 is free, and
    \texttt{start} acts while $p3$ is alive: P2 fails in state
    $(phase, captured, oldMain) = (phGap, \emptyset, \{p3\})$. When $p3$
    then binds 8421 it collides with the new process. \texttt{systemctl
    stop} returns only after the control group is empty, so Linux is safe.
    The design does not establish the same for \texttt{launchctl bootout}.
    \emph{Recommend:} capture the manager's main pid as a
    \texttt{DaemonProcess} at observe time, alongside the 8421 holders, and
    include it in condition (c). That makes P2 hold whatever the backend
    does.
\end{enumerate}

One reading the design leaves open, without a counterexample: \S4.1 does not
say whether \texttt{request\_stop} runs when the unit is already inactive but
a voxd holds 8421 (the S6 and S7 entries). The model assumes it runs, which
is idempotent for both managers. Under the resolution of F1 the answer no
longer changes any outcome.

\end{document}
